\documentclass[11pt]{article}

\usepackage[T1]{fontenc}
\usepackage{amsmath,amssymb,amsthm}
\usepackage[margin=1in]{geometry}

\newtheorem{lemma}{Lemma}

\title{Upper-Owner and Optical Cleanup}
\author{}
\date{Analytic ledger replacement, Packet E}

\begin{document}
\maketitle

\begin{abstract}
This note replaces the executable scope checks surrounding the upper owner
of the final spherical-shell residual. It proves directly that the eligible
upper wall is \(K_0(\rho)\) and that
\(k_{11}(\rho)<K_0(\rho)\) on the required ratio interval. It also records
the symbolic identities that make the repeated optical ledger samples
unnecessary.
\end{abstract}

\section{The upper owner}

For \(0<\rho<1\), set
\[
 a_\rho:=\frac{2\pi\rho}{1-\rho},
 \qquad
 \eta_\rho:=G_1\!\left(\max\{\rho,1/2\}\right),
\]
where
\[
 G_1(s)=\frac{\sqrt{1-s^2}-s\arccos s}{\pi},
 \qquad
 \omega_0:=G_1(1/2)=\frac{\sqrt3}{2\pi}-\frac16.
\]
Let \(C_0>0\), and define
\begin{equation}
 K_0(\rho):=
 \left(
  \frac{\sqrt{a_\rho}+\sqrt{a_\rho+4\eta_\rho C_0}}
       {2\eta_\rho}
 \right)^2.
 \label{eq:K0}
\end{equation}
Also put
\[
 \rho_c:=\frac1{1+2\pi},
 \qquad z_\rho:=\frac\pi{1-\rho},
 \qquad k_{11}(\rho):=\sqrt{z_\rho^2+132}.
\]

\begin{lemma}[The final upper wall is genuinely above the staircase]
\label{lem:k11-K0}
For
\[
       \rho_c\leq\rho<\frac78,
\]
one has
\begin{equation}
       k_{11}(\rho)<64<U(\rho).
       \label{eq:k11-K0}
\end{equation}
More precisely, \(U(\rho)=K_0(\rho)>64\) for
\(\rho_c\leq\rho<5/6\), while
\[
 U(\rho)=\frac{54}{(1-\rho)^2}\geq1944
 \qquad(5/6\leq\rho<7/8).
\]
In particular, on the final surviving interval
\(\rho_c\leq\rho<39/50\), the upper owner is \(K_0(\rho)\).
\end{lemma}

\begin{proof}
We use only
\[
       3<\pi<\frac{22}{7},
       \qquad \sqrt3<\frac74.
\]
They imply
\[
 \omega_0
 =\frac{\sqrt3}{2\pi}-\frac16
 <\frac{7/4}{6}-\frac16=\frac18.
\]
Since
\[
  G_1'(s)=-\frac{\arccos s}{\pi}<0\quad(0\leq s<1),
  \qquad G_1(1)=0,
\]
one has \(0<\eta_\rho\leq\omega_0<1/8\). The function
\(a_\rho\) is increasing, and its value at \(\rho_c\) is exactly one:
\[
 a_{\rho_c}
 =\frac{2\pi/(1+2\pi)}{2\pi/(1+2\pi)}=1.
\]
Thus \(a_\rho\geq1\). Formula \eqref{eq:K0} and \(C_0>0\) give
\[
 \sqrt{K_0(\rho)}
 >\frac{\sqrt{a_\rho}}{\eta_\rho}>8,
\]
so \(K_0(\rho)>64\).

 On the other hand, \(\rho<7/8\), and therefore
\[
 z_\rho<8\pi<\frac{176}{7}<26.
\]
Consequently
\[
 k_{11}(\rho)^2<26^2+132=808<64^2,
\]
which proves \eqref{eq:k11-K0}.

 The small-hole wall \(H_0\) is eligible only when
\(\rho<\omega_0\). But
\[
 \rho_c>\frac18>\omega_0,
\]
 because \(\pi<22/7<7/2\) implies \(1+2\pi<8\).
 Also \(\pi>3\) gives the exact scope chain
 \[
   \rho_c<\frac17<\frac{39}{50}<\frac56.
 \]
 For \(\rho<5/6\), the thin-seam wall is not eligible either, so the
 piecewise upper-owner formula gives \(U=K_0>64\). For
 \(5/6\leq\rho<7/8\), the seam branch owns the equality face and
 \[
  U(\rho)=\frac{54}{(1-\rho)^2}
  \geq \frac{54}{(1/6)^2}=1944>64.
 \]
 Combining this with \(k_{11}<64\) proves \eqref{eq:k11-K0}.
 Finally \(39/50<5/6\), so the final surviving interval lies wholly in
 the \(K_0\) branch.
\end{proof}

\section{Symbolic optical identities}

The optical ledger evaluates the same algebraic statements at several
sample values. The following identities replace those repetitions.

For \(a/\pi=N+t\), with \(N\geq1\) and \(0<t\leq1\), the exact
continuous-minus-product deficit is
\begin{equation}
 \frac{D(a)}{\pi^2}
 =\frac{N^2}{2}+N\left(t^2+\frac16\right)
  +\frac{2t^3}{3}.
 \label{eq:product-deficit}
\end{equation}
If \(C_D=1382/3125\), direct subtraction gives
\begin{equation}
 \frac{D(a)}{\pi^2}-C_D(N+t)^2
 =\left(\frac12-C_D\right)N^2
  +N\left(t^2-2C_Dt+\frac16\right)
  +\frac{2t^3}{3}-C_Dt^2.
 \label{eq:product-symbolic}
\end{equation}
The coefficient of \(N^2\) is positive. The difference between the
right sides at \(N+1\) and \(N\) is
\[
 2\left(\frac12-C_D\right)N
 +\left(\frac12-C_D\right)
 +t^2-2C_Dt+\frac16.
\]
The quadratic \(t^2-2C_Dt+1/6\) has minimum
\(1/6-C_D^2\) at \(t=C_D\).
Consequently, for \(N\geq1\), the displayed difference is at least
\[
 3\left(\frac12-C_D\right)+\frac16-C_D^2
 =\frac{4229603}{29296875}>0.
\]
Therefore the deficit increases in \(N\), and it is enough to check
\(N=1\). At \(N=1\), differentiation gives
\[
 \frac{d}{dt}\left(
 \frac{D(a)}{\pi^2}-C_D(1+t)^2
 \right)=2(t-C_D)(t+1).
\]
Thus the unique minimum on \(0<t\leq1\) occurs at \(t=C_D\), where the
exact value is
\begin{equation}
       \frac{1822532}{91552734375}>0.
       \label{eq:product-minimum}
\end{equation}
This proves the product screen simultaneously for every \(N,t\).

Likewise, the twelve repeated angular block checks are all instances of
\begin{equation}
       \sum_{\ell=0}^{m-1}(2\ell+1)=m^2
       \qquad(m\in\mathbb N_0),
       \label{eq:odd-sum}
\end{equation}
proved by subtracting consecutive squares. At a radial equality wall the
strict convention excludes the newly proposed layer, so no boundary
sample is added. These symbolic statements cover every repeated optical
row; the two exact endpoint reserves remain those printed in the main
optical lemma.

\section{The two optical screens with all scaling shown}

For completeness, we now print the algebra that was compressed in the main
optical lemma. Set
\[
 \varepsilon_0=\frac{11}{50},\qquad
 c=\frac{1126}{625},\qquad
 q=\frac{106}{333},\qquad
 C_D=\frac{1382}{3125}.
\]
The inequalities
\[
 \pi>\frac{333}{106},\qquad \pi<\frac{22}{7}
\]
give \(q>1/\pi\), and \(0<\varepsilon\leq\varepsilon_0<1/2\).

\begin{lemma}[Complete optical screen arithmetic]
\label{lem:optical-screens}
Let \(a=\varepsilon K\).
\begin{enumerate}
\item If \(\pi<a\leq c/\varepsilon\), then
\[
 C_D>
 \frac{2cq}{3}+\frac{\varepsilon}{4}+\varepsilon^2q^2.
 \tag{L}
\]
\item If \(a\geq c/\varepsilon\), define
\begin{align*}
 L(\varepsilon)
 &:=
 \frac{(1-\varepsilon)^2}{4}
 -\frac{(22/7)(1-\varepsilon)\varepsilon}{4c},\\
 A(\varepsilon)
 &:=
 q\varepsilon+\frac{\varepsilon^2}{4c}
 +\frac{q\varepsilon^3}{4c}
 +\frac{\varepsilon^4}{16c^2}.
\end{align*}
Then \(L(\varepsilon)>A(\varepsilon)\).
\item The low and high domains are both closed at
\(a=c/\varepsilon\), so their union covers every \(a>\pi\).
\end{enumerate}
\end{lemma}

\begin{proof}
For the low screen, \(D(a)>C_Da^2\) was proved above. The sufficient
product inequality obtained in the main optical argument is
\[
 D(a)\geq
 \varepsilon\left(\frac{2a^3}{3\pi}+\frac{a^2}{4}\right)
 +\frac{\varepsilon^2a}{\pi}.
 \tag{1}
\]
After division by \(a^2\), the three terms on the right are bounded by
\[
 \frac{2\varepsilon a}{3\pi}\leq\frac{2cq}{3},
 \qquad
 \frac{\varepsilon}{4}\leq\frac{\varepsilon_0}{4},
 \qquad
 \frac{\varepsilon^2}{\pi a}<\frac{\varepsilon^2}{\pi^2}
 <\varepsilon_0^2q^2.
\]
Exact subtraction gives
\begin{equation}
 C_D-\left(
 \frac{2cq}{3}+\frac{11}{200}
 +\left(\frac{11}{50}\right)^2q^2
 \right)
 =\frac{39569}{2772225000}>0.
 \label{eq:low-optical-reserve}
\end{equation}
This proves (L). For \(0\leq a\leq\pi\), the strict radial count is zero,
including \(a=\pi\), so no low-frequency gap remains.

For the high screen, the exact radial deficit and angular ceiling error
from the main analytic argument are
\[
 \frac{D_{\rm rad}}{a^2}
 \geq \frac{(1-\varepsilon)^2}{4}
 -\frac{\pi(1-\varepsilon)}{4a},
\]
\[
 \frac{E}{a^2}
 =\frac{\varepsilon}{\pi}
 +\frac{\varepsilon^2}{4\pi a}
 +\frac{\varepsilon}{4a}
 +\frac{\varepsilon^2}{16a^2}.
\]
Since \(a\geq c/\varepsilon\), one has \(1/a\leq\varepsilon/c\).
Using the two rational bounds for \(\pi\) therefore gives
\[
 \frac{D_{\rm rad}}{a^2}\geq L(\varepsilon),
 \qquad
 \frac{E}{a^2}<A(\varepsilon).
\]
Moreover,
\[
 L'(\varepsilon)
 =-\frac{1-\varepsilon}{2}
 -\frac{22/7}{4c}(1-2\varepsilon)<0
 \quad(0<\varepsilon<1/2),
\]
whereas \(A'(\varepsilon)>0\) coefficientwise. Hence
\[
 L(\varepsilon)-A(\varepsilon)
 \geq L(\varepsilon_0)-A(\varepsilon_0).
\]
The endpoint subtraction is
\begin{equation}
\begin{aligned}
 L(\varepsilon_0)-A(\varepsilon_0)
 ={}&
 \frac14\left(\frac{39}{50}\right)^2
 -\frac{22}{7}
  \frac{(39/50)(11/50)}{4(1126/625)}\\
 &-\left[
 \frac{11}{50}q
 +\frac{(11/50)^2}{4(1126/625)}
 +\frac{(11/50)^3q}{4(1126/625)}
 +\frac{(11/50)^4}{16(1126/625)^2}
 \right]\\
 ={}&\frac{14817541}{472867032960000}>0.
\end{aligned}
\label{eq:high-optical-reserve}
\end{equation}
Thus the radial deficit strictly exceeds the angular error.

Finally,
\[
 \frac{c}{\varepsilon}\geq\frac{c}{\varepsilon_0}
 =\frac{2252}{275}>\frac{22}{7}>\pi,
 \qquad
 \frac{2252}{275}-\frac{22}{7}
 =\frac{9714}{1925}>0.
\]
Therefore the common face lies in the low domain as well as the high
domain, and
\[
 \{a:\pi<a\leq c/\varepsilon\}
 \cup\{a:a\geq c/\varepsilon\}
 =\{a:a>\pi\},
\]
and equality \(a=c/\varepsilon\) belongs to both pieces. Together with the
zero-count interval \(0\leq a\leq\pi\), this covers every \(a\geq0\).
At \(K=a=0\), both the spectral count and the Weyl term vanish.
\end{proof}

\paragraph{Ledger rows replaced.}
Lemma~\ref{lem:k11-K0} replaces the eight executable
\texttt{verify\_u\_and\_k11} scope checks. Equations
\eqref{eq:product-deficit}--\eqref{eq:odd-sum} replace the repeated
product samples and odd-block identities. These are rows of
\texttt{verify\_optical\_constants}.
Lemma~\ref{lem:optical-screens}
prints the remaining scaling, monotonicity, common-face, and exact-reserve
implications. Thus all fifty optical rows are consequences of displayed
identities rather than executable predicates.

\end{document}
